\documentclass[10pt,a4paper]{amsart}
\usepackage[margin=19mm]{geometry}
\usepackage{amsmath,amssymb,mathtools}
\usepackage[scr=rsfs,cal=euler]{mathalfa}
\usepackage{xfrac}
\usepackage[unicode,hidelinks,bookmarks=false]{hyperref}
\newcommand{\ie}{i.e.\ }
\newcommand{\cf}{cf.\ }
\newcommand{\resp}{respectively\ }
\newcommand{\Subsection}{\subsection}
\providecommand{\nobreakdash}{\nobreak\hbox{-}}
\setlength{\emergencystretch}{2em}
\multlinegap0pt
\usepackage{bm}
\usepackage{bbm}
\usepackage{dsfont}
\usepackage{xspace}
\usepackage{cancel}
\usepackage{mathtools}
\usepackage{amsthm}
\newtheorem{theorem}{Theorem}
\newtheorem{proposition}[theorem]{Proposition}
\newcommand{\Zm}{\mathbb Z_m}
\newcommand{\cO}{\Omega}
\title[Hamilton decompositions of equal-side directed tori]{Hamilton decompositions of equal-side directed tori (Proposition 4)}
\author{SangHyun Park}
\date{Source: arXiv:2605.04734v3 (CC BY 4.0)}
\begin{document}
\maketitle

\noindent\emph{Source and licence.} Hamilton decompositions of equal-side directed tori. Version arXiv:2605.04734v3 (CC BY 4.0), distributed under the Creative Commons Attribution 4.0 International licence, as recorded in the arXiv licence metadata for this exact version and shown on the abstract page https://arxiv.org/abs/2605.04734v3. Full rights notice: licenses/arxiv-2605.04734-CC-BY-4.0.txt.\par\smallskip

\noindent\textit{Editorial scope.} Counted unit: the Proposition printed as Proposition 4 of the article (source0.tex lines 582--595, source label \texttt{prop:pinned-selection}), delivered together with the article's own complete original proof of it (source0.tex lines 596--662, one \texttt{proof} environment of 67 lines). No mathematics is written, repaired or re-derived: statement and proof are the article's own text line for line. Required context retained uncounted: none. Every definition, display and label of those lines is retained unchanged. Omitted from this package and not redistributed: 1--581, 663--3562. Nothing inside a retained excerpt is omitted or altered: every retained body is delivered exactly as the article's lines are written, so no omission receipt of any kind accompanies this package. Citations: the retained excerpts carry no citation command at all, so no bibliography entry of the article is transcribed and no reference is redistributed. Reference adaptation: none was needed and none was made; every reference inside the delivered unit resolves inside it, so no unresolved reference or citation is produced. Printed slips kept literal: no printed word, symbol, hypothesis or number of the counted statement or of its proof was altered. Layout and class: the article's own class is \texttt{amsart} with packages including fontenc, lmodern, microtype, geometry, amsmath, amssymb, amsthm, mathtools, booktabs, array, hyperref, enumitem; the wrapper is the lane's pinned \texttt{amsart} editorial wrapper at 10pt on A4, which restates the theorem-like environments the retained text needs and adds the article's own macro definitions that the retained text needs (\texttt{\textbackslash Zm}, \texttt{\textbackslash cO}), with extra wrapper packages (bm, bbm, dsfont, xspace, cancel, mathtools). The delivered numbering is the unit's own. Assets: no retained span contains a figure environment, a graphics-inclusion command, a PDF-page inclusion, a table or a TikZ picture; no image file is copied into this package, the package declares no asset dependency and no image file is redistributed with it. Licence: version arXiv:2605.04734v3 of this article is distributed under the Creative Commons Attribution 4.0 International licence, as recorded in the arXiv licence metadata for this exact version and shown on the abstract page https://arxiv.org/abs/2605.04734v3. A full rights notice is delivered at licenses/arxiv-2605.04734-CC-BY-4.0.txt. Characters: every retained excerpt is pure ASCII, so the delivered engine, XeLaTeX with its default Latin Modern text font, reports no missing character.
\begin{proposition}[Connected selection with prescribed exclusions]\label{prop:pinned-selection}
Let $m\ge4$ be even, and let $\cO_\mathrm{dist}\subseteq\cO$ be a set of
distinguished columns satisfying
$|A_B\cap\cO_\mathrm{dist}|\le1$ for every $B$. The following are equivalent.
\begin{enumerate}
\item Each component of the incidence graph contains an even number of column
vertices.
\item One can choose $J_{B,t}\subseteq A_B$, with $|J_{B,t}|=b_B$, so that every
column total is $1$ or $-1$ modulo $m$, no distinguished column belongs to $J_{B,0}$,
and, for each $B$, both families
$\{J_{B,t}:t\in\Zm\}$ and $\{A_B\setminus J_{B,t}:t\in\Zm\}$ are
intersection-connected whenever their members have size at least two.
\end{enumerate}
\end{proposition}

\begin{proof}
For necessity, sum the selected entries in an incidence component. Each block
contributes $mb_B$, an even number, while every column contributes an odd number.

For sufficiency, we first choose and orient pairs of columns, then realize
each oriented pair by a row selection. The choice of exceptional rows will
provide the two connectivity conclusions.

\emph{Parity and orientation.}
Take a spanning tree in each component of the incidence graph. Select its
edges so that column vertices have odd selected degree and block vertices
have even selected degree. To construct this selection, remove leaves
successively, selecting a leaf's parent edge exactly when its remaining
parity requirement is odd, and updating the requirement at its parent.
The last requirement is zero because the component has an even number of
column vertices.
Pair the selected incidences at each block. Since each column occurs at most
once in that block, these are disjoint local pairs. Across all blocks, they
define a multigraph on $\cO$ in which every vertex has odd degree; retain the
block label on each edge.

For each connected component of this multigraph, introduce one new vertex
and join it once to every vertex of that component. The component has an
even number of vertices by the handshaking lemma, so the new vertex and all
old vertices have even degree. Orient an Euler circuit of the enlarged
component and remove the new edges. At every column, the remaining indegree
minus outdegree is $+1$ or $-1$. This orients all local pairs while retaining
their block labels.

\emph{Row quotas and column totals.}
In a block with $a$ columns, write $s$ for the number of local pairs and put
$b=\lfloor a/2\rfloor$. For a directed pair $u\to v$, choose an exceptional
row $\eta$: select $v$ on row $\eta$ and $u$ on each of the other $m-1$ rows.
The pair contributes exactly one selected entry to every row and contributes
$-1,+1$ modulo $m$ at $u,v$, respectively. Thus the column totals of all
pairs equal the oriented divergences modulo $m$.
Choose a further $b-s$ unpaired columns on every row. If the distinguished
column is unpaired, exclude it from this fixed choice. There are enough
columns: for $a=2b$, an unpaired distinguished column gives $s\le b-1$ and
\[
 a-2s-1\ge b-s;
\]
for $a=2b+1$, the same inequality follows from $s\le b$. A paired
distinguished column is already outside the unpaired set. The resulting
rows have size $b$, and every fixed selected column contributes zero modulo
$m$.

\emph{Exclusions and connectivity.}
If the distinguished column is the first endpoint of a pair, put that
pair's exceptional row at $0$; if it is the second endpoint, put the
exceptional row at $1$. When $s\ge2$, assign a second pair to the other of
rows $0,1$, and any remaining pairs to the first of these rows. With no
distinguished pair, use row $0$ for a sole pair; for $s\ge2$, assign the
first two pairs to rows $0,1$ and the rest to row $0$. If $s=0$, all rows
are constant. The distinguished column is excluded from row $0$ in each case.

Row $2$ contains the first endpoint of every pair. A fixed selected column
connects all selected supports. If there is no such column and $b\ge2$,
then $s=b\ge2$. Each exceptional row shares an endpoint with row $2$
through a pair whose exceptional row is the other of $0,1$. All other rows
agree with row $2$, proving intersection connectivity.
For complements, a fixed unselected column again connects all rows.
Otherwise every unpaired column was selected, so the complementary size
is $s$; if this size is at least two, the same argument applies to the
second endpoints on row $2$. Both row families therefore have the required
connectivity, and the column totals are the divergences $\pm1$ modulo $m$.
\end{proof}

\end{document}
